\documentclass[aspectratio=169, 10pt]{beamer}

\usetheme{Madrid}
\usefonttheme{professionalfonts}

\usepackage[utf8]{inputenc}
\usepackage[indonesian]{babel}
\usepackage{amsmath,amsfonts,amssymb}
\usepackage{graphicx}
\usepackage{booktabs}
\usepackage{colortbl}
\usepackage{listings}
\usepackage{xcolor}

\definecolor{mainblue}{RGB}{24, 43, 73}
\definecolor{accentblue}{RGB}{41, 128, 185}
\definecolor{accentgold}{RGB}{230, 126, 34}
\definecolor{darkgray}{RGB}{44, 62, 80}
\definecolor{lightbg}{RGB}{248, 250, 252}
\definecolor{bordergray}{RGB}{200, 210, 220}
\definecolor{codebg}{RGB}{240, 243, 246}
\definecolor{codegreen}{RGB}{39, 174, 96}

\setbeamercolor{structure}{fg=mainblue}
\setbeamercolor{palette primary}{bg=mainblue, fg=white}
\setbeamercolor{palette secondary}{bg=accentblue, fg=white}
\setbeamercolor{palette tertiary}{bg=mainblue!80, fg=white}
\setbeamercolor{title}{fg=white, bg=mainblue}
\setbeamercolor{frametitle}{fg=white, bg=mainblue}
\setbeamercolor{block title}{bg=mainblue, fg=white}
\setbeamercolor{block body}{bg=lightbg, fg=darkgray}
\setbeamercolor{block title alerted}{bg=accentgold, fg=white}
\setbeamercolor{block body alerted}{bg=accentgold!10, fg=darkgray}

\lstset{
    language=Python,
    backgroundcolor=\color{codebg},
    basicstyle=\ttfamily\scriptsize\color{darkgray},
    keywordstyle=\bfseries\color{accentblue},
    stringstyle=\color{accentgold},
    commentstyle=\itshape\color{codegreen},
    showstringspaces=false,
    breaklines=true,
    frame=single,
    rulecolor=\color{bordergray},
    numbers=none
}

\title[\textbf{Deret Waktu (Minggu 10)}]{\textbf{Feature Engineering \& Machine Learning untuk STLF}}
\subtitle{Teknik Sliding Window, Fitur Lag/Kalender, XGBoost, dan Validasi Silang Temporal}
\author[\textbf{Ir. Novalio Daratha, Ph.D.}]{\textbf{Ir. Novalio Daratha, S.T., M.Sc., Ph.D.}}
\institute[FT UNIB]{Program Studi Teknik Elektro -- Fakultas Teknik\\Universitas Bengkulu}
\date{Minggu 10 / Semester Ganjil 2026}

\begin{document}

\begin{frame}[plain]
    \titlepage
\end{frame}

\begin{frame}{Agenda Perkuliahan Minggu ke-10}
    \begin{columns}[T]
        \begin{column}{0.48\textwidth}
            \begin{block}{Capaian Pembelajaran (Sub-CPMK 9)}
                Mahasiswa mampu merumuskan problem peramalan deret waktu ke dalam kerangka \textit{Supervised Learning}, merancang rekayasa fitur temporal, serta melatih dan memvalidasi model \textit{Gradient Boosting} (XGBoost/LightGBM).
            \end{block}
        \end{column}
        \begin{column}{0.48\textwidth}
            \begin{enumerate}
                \item \textbf{Formulasi Supervised Learning}
                \item \textbf{Teknik Rekayasa Fitur Temporal}
                \item \textbf{Encoding Kalender Siklis (Sin/Cos)}
                \item \textbf{Algoritma Tree-Based: XGBoost \& LightGBM}
                \item \textbf{Time Series Cross-Validation (Mencegah Data Leakage)}
                \item \textbf{Feature Importance \& Analisis Hasil}
            \end{enumerate}
        \end{column}
    \end{columns}
\end{frame}

\begin{frame}{Mengubah Deret Waktu Menjadi Supervised Learning}
    Menggunakan teknik \textbf{Sliding Window / Lagged Matrix}:
    \begin{block}{Representasi Matriks Fitur ($X$) dan Target ($y$)}
        \begin{equation}
            \begin{bmatrix}
                y_t \\
                y_{t+1} \\
                \vdots \\
                y_{N}
            \end{bmatrix}
            = f\left(
            \begin{bmatrix}
                y_{t-1} & y_{t-2} & \dots & y_{t-p} & X_{\text{kalender}, t} & X_{\text{cuaca}, t} \\
                y_t & y_{t-1} & \dots & y_{t-p+1} & X_{\text{kalender}, t+1} & X_{\text{cuaca}, t+1} \\
                \vdots & \vdots & \ddots & \vdots & \vdots & \vdots \\
                y_{N-1} & y_{N-2} & \dots & y_{N-p} & X_{\text{kalender}, N} & X_{\text{cuaca}, N}
            \end{bmatrix}
            \right)
        \end{equation}
    \end{block}
    Keunggulan: Dapat menyerap puluhan fitur non-linier secara simultan tanpa asumsi distribusi kaku.
\end{frame}

\begin{frame}{Taksonomi Rekayasa Fitur (*Feature Engineering*)}
    \begin{columns}[T]
        \begin{column}{0.33\textwidth}
            \begin{block}{1. Fitur Lag Historis}
                \begin{itemize}
                    \item Lag dekat: $t-1, t-2, t-3$
                    \item Lag harian: $t-24, t-48$
                    \item Lag mingguan: $t-168$
                \end{itemize}
            \end{block}
        \end{column}
        \begin{column}{0.33\textwidth}
            \begin{block}{2. Statistik Bergerak}
                \begin{itemize}
                    \item Rolling Mean (24 jam)
                    \item Rolling Std Dev
                    \item Rolling Max / Min
                    \item Exponential Weighted MA
                \end{itemize}
            \end{block}
        \end{column}
        \begin{column}{0.33\textwidth}
            \begin{block}{3. Fitur Kalender}
                \begin{itemize}
                    \item Jam dalam sehari (0--23)
                    \item Hari dalam seminggu (0--6)
                    \item Hari Libur Nasional (Flag biner 0/1)
                \end{itemize}
            \end{block}
        \end{column}
    \end{columns}
    \vspace{0.2cm}
    \begin{alertblock}{Encoding Siklik (Jam \& Bulan)}
        Nilai jam 23 dan jam 00 berdekatan secara fisik. Gunakan transformasi trigonometri:
        $x_{\sin} = \sin(2\pi \cdot \text{jam} / 24)$ dan $x_{\cos} = \cos(2\pi \cdot \text{jam} / 24)$.
    \end{alertblock}
\end{frame}

\begin{frame}{Mencegah Data Leakage: Time Series Cross-Validation}
    \begin{columns}[T]
        \begin{column}{0.48\textwidth}
            \begin{block}{Bahaya $K$-Fold CV Acak Standar}
                Pengacakan baris (*shuffling*) menyebabkan data masa depan bocor ke masa lalu (\textit{Lookahead Bias / Leakage}). Model tampak memiliki akurasi $99\%$ di lab namun gagal total di lapangan.
            \end{block}
        \end{column}
        \begin{column}{0.48\textwidth}
            \begin{block}{Metode Validasi yang Benar}
                \textbf{Expanding Window (Time Series Split)}:
                \begin{itemize}
                    \item Fold 1: Train [Bulan 1--3] $\to$ Test [Bulan 4]
                    \item Fold 2: Train [Bulan 1--4] $\to$ Test [Bulan 5]
                    \item Fold 3: Train [Bulan 1--5] $\to$ Test [Bulan 6]
                \end{itemize}
            \end{block}
        \end{column}
    \end{columns}
\end{frame}

\begin{frame}[fragile]{Implementasi Python: XGBoost untuk STLF}
\begin{lstlisting}
import xgboost as xgb
from sklearn.metrics import mean_squared_error

# 1. Pipeline Feature Engineering
df['hour_sin'] = np.sin(2 * np.pi * df.index.hour / 24.0)
df['hour_cos'] = np.cos(2 * np.pi * df.index.hour / 24.0)
df['lag_24'] = df['load'].shift(24)
df['roll_mean_24'] = df['load'].shift(1).rolling(24).mean()
data = df.dropna()

# 2. Split Berurutan Temporal
X = data[['hour_sin', 'hour_cos', 'lag_24', 'roll_mean_24', 'temp']]
y = data['load']
X_train, X_test = X[:-168], X[-168:]
y_train, y_test = y[:-168], y[-168:]

# 3. Training Gradient Boosting
model = xgb.XGBRegressor(n_estimators=1000, learning_rate=0.02, max_depth=5, subsample=0.8)
model.fit(X_train, y_train, eval_set=[(X_test, y_test)], early_stopping_rounds=50, verbose=False)

# 4. Feature Importance
xgb.plot_importance(model, max_num_features=5)
\end{lstlisting}
\end{frame}

\begin{frame}{Rangkuman \& Persiapan Minggu ke-11}
    \begin{block}{Poin Kunci Minggu 10}
        \begin{itemize}
            \item Rekayasa fitur (lag, rolling, calendar encoding) adalah kunci akurasi algoritma Machine Learning.
            \item XGBoost dan LightGBM sangat tangguh menangani relasi non-linier kompleks beban sistem tenaga.
            \item Validasi temporal yang ketat tanpa shuffling menjamin keandalan model saat diterapkan pada sistem online.
        \end{itemize}
    \end{block}
    \vspace{0.2cm}
    \begin{alertblock}{Materi Minggu Depan (Minggu ke-11)}
        \textbf{Deep Learning Dasar untuk Deret Waktu}: Arsitektur Recurrent Neural Network (RNN), Long Short-Term Memory (LSTM), dan Gated Recurrent Unit (GRU) untuk data pembangkitan EBT.
    \end{alertblock}
\end{frame}

\begin{frame}[plain]
    \begin{center}
        \vspace{1.5cm}
        {\Huge \color{mainblue} \textbf{Terima Kasih}}\\[0.4cm]
        {\large Pertanyaan \& Diskusi Machine Learning untuk Deret Waktu}\\[1.2cm]
        \textbf{Laboratorium Sistem Tenaga \& Komputasi Cerdas}\\
        Jurusan Teknik Elektro -- Universitas Bengkulu
    \end{center}
\end{frame}

\end{document}
